\documentclass[11pt,a4paper]{article}
\usepackage[a4paper,margin=1in,headsep=30pt,footskip=35pt]{geometry}
\usepackage[T1]{fontenc}
\usepackage[utf8]{inputenc}
\usepackage{tgpagella}
\usepackage[scale=0.92]{tgheros}
\usepackage{microtype}
\usepackage{graphicx}
\usepackage{hyperref}
\usepackage{enumitem}
\usepackage{xcolor}
\usepackage{titlesec}
\usepackage{parskip}
\usepackage{longtable}
\usepackage{booktabs}
\usepackage{array}
\usepackage{tabularx}
\usepackage{fancyhdr}
\usepackage{lastpage}
\usepackage{tikz}
\usepackage[most]{tcolorbox}

% --- CORE COLOR IDENTITY ---
\definecolor{KazePrimary}{HTML}{284F58}
\definecolor{KazeSecondary}{HTML}{8A6538}
\definecolor{KazeInk}{HTML}{1F292D}
\definecolor{KazeMuted}{HTML}{5D676C}
\definecolor{KazeSoft}{HTML}{F8F4ED}
\definecolor{KazeLine}{HTML}{D8CCBC}
\definecolor{KazePaper}{HTML}{FCFBF8}
\definecolor{KazeLegal}{HTML}{7A2E2E}

\hypersetup{colorlinks=true, linkcolor=KazePrimary, urlcolor=KazePrimary}

\setlist[itemize]{leftmargin=1.5em, itemsep=0.4em}
\renewcommand{\arraystretch}{1.25}
\setlength{\headheight}{35pt}

\newcommand{\KazeLogoPath}{../../composeApp/src/commonMain/composeResources/drawable/k_logo_raster.png}
\newcommand{\KazeMarkPath}{../../composeApp/src/commonMain/composeResources/drawable/k_mark_raster.png}
\newcommand{\GaboLogoPath}{../../composeApp/src/commonMain/composeResources/drawable/gabo_mark_raster.png}
\newcommand{\GaboGrayMarkPath}{../../composeApp/src/commonMain/composeResources/drawable/gabo_mark_gray_raster.png}
\newcommand{\KazeWordmarkHeader}{{\fontfamily{qhv}\selectfont\sffamily\bfseries\small\textcolor{KazePrimary}{KAZE}}}
\newcommand{\KazeWordmarkCover}{{\fontfamily{qhv}\selectfont\sffamily\bfseries\fontsize{30}{30}\selectfont\textcolor{KazePrimary}{KAZE}}}
\newcommand{\KazeByGaboHeader}{{\raisebox{0.05cm}{\includegraphics[height=0.18cm]{\GaboLogoPath}}\hspace{0.2em}\sffamily\tiny\textcolor{KazeSecondary}{by GABO}}}
\newcommand{\KazeByGaboCover}{{\includegraphics[height=0.45cm]{\GaboLogoPath}\hspace{0.4em}\sffamily\large\textcolor{KazeSecondary}{by GABO}}}
\newcommand{\KazeCoverLabel}[1]{{\sffamily\normalsize\textcolor{KazeSecondary}{\textbf{#1}}}}
\newcommand{\KazeCoverTitle}[1]{{\fontfamily{qpl}\selectfont\fontsize{38}{43}\selectfont\textcolor{KazeInk}{#1}}}
\newcommand{\KazeCoverSubtitle}[1]{{\sffamily\Large\textcolor{KazePrimary}{#1}}}
\newcommand{\KazeContactEmail}{orestegabo@icloud.com}
\newcommand{\KazeContactWebsiteUrl}{https://kazerwanda.com}
\newcommand{\KazeContactWebsiteLabel}{kazerwanda.com}
\newcommand{\KazeContactBlock}{
    \begin{itemize}
        \item \textbf{Email:} \href{mailto:\KazeContactEmail}{\KazeContactEmail}
        \item \textbf{Website:} \href{\KazeContactWebsiteUrl}{\KazeContactWebsiteLabel}
    \end{itemize}
}

\pagecolor{KazePaper}

\AddToHook{shipout/background}{%
    \begin{tikzpicture}[remember picture,overlay]
        \fill[KazePrimary, opacity=0.018, rounded corners=14pt]
            ([xshift=0.04\paperwidth,yshift=-0.07\paperheight]current page.north west)
            rectangle
            ([xshift=0.96\paperwidth,yshift=-0.31\paperheight]current page.north west);
        \fill[KazeSecondary, opacity=0.014, rounded corners=14pt]
            ([xshift=0.05\paperwidth,yshift=0.34\paperheight]current page.north west)
            rectangle
            ([xshift=0.95\paperwidth,yshift=0.14\paperheight]current page.south west);

        \draw[KazePrimary, line width=1.7pt, opacity=0.09, line cap=round]
            ([xshift=0.08\paperwidth,yshift=-0.10\paperheight]current page.north west) --
            ([xshift=0.72\paperwidth,yshift=-0.10\paperheight]current page.north west);
        \draw[KazePrimary, line width=0.95pt, opacity=0.085, line cap=round]
            ([xshift=0.12\paperwidth,yshift=-0.135\paperheight]current page.north west) --
            ([xshift=0.88\paperwidth,yshift=-0.135\paperheight]current page.north west);
        \draw[KazeSecondary, line width=1.2pt, opacity=0.08, line cap=round]
            ([xshift=0.18\paperwidth,yshift=0.26\paperheight]current page.south west) --
            ([xshift=0.84\paperwidth,yshift=0.26\paperheight]current page.south west);
        \draw[KazePrimary, line width=1.7pt, opacity=0.09, line cap=round]
            ([xshift=0.14\paperwidth,yshift=0.21\paperheight]current page.south west) --
            ([xshift=0.62\paperwidth,yshift=0.21\paperheight]current page.south west);

        \draw[KazePrimary, line width=1.15pt, opacity=0.09]
            ([xshift=0.62\paperwidth,yshift=-0.06\paperheight]current page.north west) --
            ([xshift=0.76\paperwidth,yshift=-0.06\paperheight]current page.north west) --
            ([xshift=0.82\paperwidth,yshift=-0.11\paperheight]current page.north west) --
            ([xshift=0.93\paperwidth,yshift=-0.11\paperheight]current page.north west) --
            ([xshift=0.93\paperwidth,yshift=-0.18\paperheight]current page.north west) --
            ([xshift=0.82\paperwidth,yshift=-0.18\paperheight]current page.north west) --
            ([xshift=0.75\paperwidth,yshift=-0.24\paperheight]current page.north west) --
            ([xshift=0.58\paperwidth,yshift=-0.24\paperheight]current page.north west);

        \draw[KazeSecondary, line width=1.15pt, opacity=0.08]
            ([xshift=0.10\paperwidth,yshift=0.12\paperheight]current page.south west) --
            ([xshift=0.26\paperwidth,yshift=0.12\paperheight]current page.south west) --
            ([xshift=0.33\paperwidth,yshift=0.17\paperheight]current page.south west) --
            ([xshift=0.48\paperwidth,yshift=0.17\paperheight]current page.south west) --
            ([xshift=0.48\paperwidth,yshift=0.10\paperheight]current page.south west) --
            ([xshift=0.34\paperwidth,yshift=0.10\paperheight]current page.south west) --
            ([xshift=0.27\paperwidth,yshift=0.05\paperheight]current page.south west) --
            ([xshift=0.12\paperwidth,yshift=0.05\paperheight]current page.south west);

        \draw[KazePrimary, line width=1.8pt, opacity=0.05]
            ([xshift=0.70\paperwidth,yshift=-0.22\paperheight]current page.north west) circle [radius=2.0cm];
        \draw[KazeSecondary, line width=1.8pt, opacity=0.048]
            ([xshift=0.14\paperwidth,yshift=0.18\paperheight]current page.south west) circle [radius=2.35cm];

        \draw[KazePrimary, line cap=round, opacity=0.065, line width=1.2pt]
            ([xshift=0.11\paperwidth,yshift=-0.185\paperheight]current page.north west) --
            ([xshift=0.36\paperwidth,yshift=-0.185\paperheight]current page.north west);
        \draw[KazeSecondary, line cap=round, opacity=0.06, line width=0.8pt]
            ([xshift=0.11\paperwidth,yshift=-0.203\paperheight]current page.north west) --
            ([xshift=0.42\paperwidth,yshift=-0.203\paperheight]current page.north west);
        \draw[KazePrimary, line cap=round, opacity=0.065, line width=1.2pt]
            ([xshift=0.11\paperwidth,yshift=-0.221\paperheight]current page.north west) --
            ([xshift=0.48\paperwidth,yshift=-0.221\paperheight]current page.north west);
        \draw[KazeSecondary, line cap=round, opacity=0.06, line width=0.8pt]
            ([xshift=0.11\paperwidth,yshift=-0.239\paperheight]current page.north west) --
            ([xshift=0.54\paperwidth,yshift=-0.239\paperheight]current page.north west);
        \draw[KazePrimary, line cap=round, opacity=0.065, line width=1.2pt]
            ([xshift=0.11\paperwidth,yshift=-0.257\paperheight]current page.north west) --
            ([xshift=0.60\paperwidth,yshift=-0.257\paperheight]current page.north west);
        \draw[KazeSecondary, line cap=round, opacity=0.06, line width=0.8pt]
            ([xshift=0.11\paperwidth,yshift=-0.275\paperheight]current page.north west) --
            ([xshift=0.66\paperwidth,yshift=-0.275\paperheight]current page.north west);
        \draw[KazePrimary, line cap=round, opacity=0.065, line width=1.2pt]
            ([xshift=0.11\paperwidth,yshift=-0.293\paperheight]current page.north west) --
            ([xshift=0.72\paperwidth,yshift=-0.293\paperheight]current page.north west);
        \node[anchor=south east, opacity=0.035] at ([xshift=-0.45cm,yshift=0.65cm]current page.south east) {%
            \includegraphics[width=4.8cm]{\KazeMarkPath}%
        };
        \node[anchor=north east, opacity=0.028] at ([xshift=-1.1cm,yshift=-1.2cm]current page.north east) {%
            \includegraphics[width=2.2cm]{\GaboGrayMarkPath}%
        };
    \end{tikzpicture}%
}

% --- SECTION DESIGN: TECHNICAL PRECISION ---
\titleformat{\section}
{\sffamily\Large\bfseries\color{KazeInk}}
{\thesection}
{1em}
{\MakeUppercase}
[\vspace{0.2em}{\color{KazeSecondary}\hrule height 1.5pt width 2cm}\vspace{0.5em}]

\titleformat{\subsection}
{\sffamily\large\bfseries\color{KazePrimary}}
{\thesubsection}
{1em}
{}

% --- DUAL-BRANDED HEADER (HIGH VISIBILITY) ---
\pagestyle{fancy}
\fancyhf{}
\renewcommand{\headrulewidth}{0pt}
\renewcommand{\footrulewidth}{0pt}

\fancyhead[L]{%
    \begin{tabular}{@{}l l@{}}
        \raisebox{-0.3\height}{\includegraphics[height=0.6cm]{\KazeLogoPath}} &
        \begin{minipage}[b]{4cm}
            \KazeWordmarkHeader\\
            \vspace{-1mm}
            \KazeByGaboHeader
        \end{minipage}
    \end{tabular}
}
\fancyhead[R]{\sffamily\tiny\textcolor{KazeMuted}{KAZE PROPRIETARY // \today}}

\fancyfoot[L]{%
    \begin{tabular}{@{}l@{}}
        {\color{KazeLine}\rule{3.1cm}{0.6pt}}\\[-0.2em]
        {\sffamily\tiny\textcolor{KazeLegal}{CONFIDENTIAL PRIVATE DOCUMENT}}
    \end{tabular}
}
\fancyfoot[C]{%
    \sffamily\scriptsize
    \textcolor{KazeMuted}{\textsc{Page} \thepage\ of \pageref*{LastPage}}
}
\fancyfoot[R]{%
    \begin{tabular}{@{}r@{}}
        {\color{KazeLine}\rule{2.1cm}{0.6pt}}\\[-0.2em]
        {\sffamily\tiny\textcolor{KazePrimary}{Kaze by GABO}}
    \end{tabular}
}

% --- CLEAN REGISTRY BOXES ---
\newtcolorbox{kazehighlightbox}{
    enhanced,
    colback=KazeSoft,      % Only a very faint paper-like tint
    frame hidden,          % No borders at all
    sharp corners,
    left=5mm, right=5mm, top=4mm, bottom=4mm,
% Use a small square bullet as an anchor instead of a line
    overlay={
        \fill[KazeSecondary] ([xshift=2mm,yshift=-4mm]frame.north west) rectangle ([xshift=3mm,yshift=-5mm]frame.north west);
    }
}
\newtcolorbox{kazecalloutbox}[1][]{
    enhanced,
    colback=white,
    colframe=KazeInk,
    boxrule=1pt,
    sharp corners,
    fonttitle=\sffamily\bfseries\small,
    coltitle=white,
    attach boxed title to top left={yshift=-2mm, xshift=5mm},
    boxed title style={colback=KazeInk, sharp corners, boxrule=0pt},
    title=#1,
    left=5mm, top=6mm, bottom=4mm
}

\newcommand{\KazePageTitle}[1]{%
    \vspace{1em}
    {\sffamily\small\textcolor{KazeSecondary}{\textbf{\MakeUppercase{Document Scope: #1}}}\par}
    \vspace{0.5em}
}

\newcommand{\KazeDocumentNotice}{%
    \begin{kazecalloutbox}[Security Disclosure]
    \sffamily\small This dossier contains trade secrets and sensitive intellectual property. \textbf{Unauthorized distribution is strictly prohibited.} Access is limited to Kaze personnel and authorized external stakeholders for legal and investor review.
    \end{kazecalloutbox}
}

% --- THE "Blueprints" COVER PAGE ---
\newcommand{\KazeCover}[4]{
    \hypersetup{pageanchor=false}
    \begin{titlepage}
    \thispagestyle{empty}

% Background "Pillar" for institutional weight
    \begin{tikzpicture}[remember picture,overlay]
    \fill[KazePrimary] (current page.north west) rectangle ([xshift=0.5cm]current page.south west);
    \draw[KazePrimary!55, line width=1.8pt, opacity=0.22, line cap=round]
        ([xshift=0.08\paperwidth,yshift=-0.10\paperheight]current page.north west) --
        ([xshift=0.72\paperwidth,yshift=-0.10\paperheight]current page.north west);
    \draw[KazeSecondary!70, line width=1.1pt, opacity=0.19, line cap=round]
        ([xshift=0.12\paperwidth,yshift=-0.135\paperheight]current page.north west) --
        ([xshift=0.88\paperwidth,yshift=-0.135\paperheight]current page.north west);
    \draw[KazePrimary!55, line width=1.2pt, opacity=0.20]
        ([xshift=0.62\paperwidth,yshift=-0.06\paperheight]current page.north west) --
        ([xshift=0.76\paperwidth,yshift=-0.06\paperheight]current page.north west) --
        ([xshift=0.82\paperwidth,yshift=-0.11\paperheight]current page.north west) --
        ([xshift=0.93\paperwidth,yshift=-0.11\paperheight]current page.north west) --
        ([xshift=0.93\paperwidth,yshift=-0.18\paperheight]current page.north west) --
        ([xshift=0.82\paperwidth,yshift=-0.18\paperheight]current page.north west) --
        ([xshift=0.75\paperwidth,yshift=-0.24\paperheight]current page.north west) --
        ([xshift=0.58\paperwidth,yshift=-0.24\paperheight]current page.north west);
    \draw[KazeSecondary!75, line width=1.35pt, opacity=0.18, line cap=round]
        ([xshift=0.18\paperwidth,yshift=0.26\paperheight]current page.south west) --
        ([xshift=0.84\paperwidth,yshift=0.26\paperheight]current page.south west);
    \draw[KazePrimary!55, line width=1.8pt, opacity=0.20, line cap=round]
        ([xshift=0.14\paperwidth,yshift=0.21\paperheight]current page.south west) --
        ([xshift=0.62\paperwidth,yshift=0.21\paperheight]current page.south west);
    \end{tikzpicture}

    \vspace*{0.5cm}
    \begin{flushleft}
    \hspace{1cm}
    \begin{tabular}{@{}l l@{}}
    \raisebox{-0.3\height}{\includegraphics[width=2.2cm]{\KazeLogoPath}} &
    \begin{minipage}[b]{8cm}
    {\KazeWordmarkCover\par}
    \vspace{0.1cm}
    {\KazeByGaboCover\par}
    \end{minipage}
    \end{tabular}
    \end{flushleft}

    \vspace{3.5cm}

    \hspace{1cm}
    \begin{minipage}{0.85\textwidth}
    \raggedright
    {\KazeCoverLabel{OFFICIAL DOSSIER / #4}\par}
    \vspace{0.95cm}
    {\KazeCoverTitle{#1}\par}
    \vspace{0.7cm}
    {\KazeCoverSubtitle{#2}\par}

    \vspace{1.5cm}
    {\color{KazeLine}\hrule height 2pt width 3cm}
    \vspace{1.5cm}

    {\sffamily\large\textcolor{KazeMuted}{#3}\par}
    \end{minipage}

    \vfill

    \hspace{1cm}
    \begin{minipage}{0.85\textwidth}
    \begin{kazehighlightbox}
    \sffamily\small
    \textbf{\textcolor{KazeLegal}{RESTRICTED DOCUMENT:}} This information is intended for the recipient only. It contains proprietary strategic planning belonging to Kaze. By accessing this document, you agree to non-disclosure terms.
    \end{kazehighlightbox}
    \end{minipage}

    \vspace{1.5cm}
    \hspace{1cm}
    \begin{minipage}{0.85\textwidth}
    \sffamily\tiny\textcolor{KazeMuted}{AUTHORIZATION \& CONTACT:} \\
    \sffamily\small\textbf{\href{mailto:\KazeContactEmail}{\KazeContactEmail}} \hfill \textbf{\MakeUppercase{\KazeContactWebsiteLabel}}
    \end{minipage}

    \end{titlepage}
    \hypersetup{pageanchor=true}
}


\begin{document}

\KazeCover
  {Kaze Cahier des Charges}
  {Requirements specification for the Kaze mobile platform, venue marketplace, guest access, and partner ecosystem}
  {April 2026}
  {Product and technical planning}

\KazePageTitle{Cahier des Charges}
\KazeDocumentNotice

\section{Document Purpose}

This document is the formal \textit{cahier des charges} for Kaze. In English, it can be read as a requirements specification or product and technical specification.

It describes what Kaze is expected to support, who it serves, which workflows matter first, and which constraints should guide future product and engineering decisions.

\begin{kazecalloutbox}[Intended Use]
This document can be used for internal planning, future developer onboarding, investor discussions, partner conversations, and technical scoping before implementation work begins.
\end{kazecalloutbox}

\section{Project Identity}

\begin{longtable}{p{0.28\textwidth} p{0.64\textwidth}}
\textbf{Topic} & \textbf{Requirement} \\
\midrule
Product name & Kaze Rwanda / Kaze by GABO. \\
Product type & Mobile-first hospitality, venue discovery, invitation, access-pass, and service-request platform. \\
Primary market & Rwanda, with an initial focus on hotels, wedding venues, conference rooms, local stays, and event service providers. \\
Business status & Private commercial product. Not open source. Use, copying, contribution, or distribution require authorization from GABO or a separate agreement. \\
Primary contact & \href{mailto:\KazeContactEmail}{\KazeContactEmail} and \href{\KazeContactWebsiteUrl}{\KazeContactWebsiteLabel}. \\
\end{longtable}

\section{Product Vision}

Kaze should reduce the friction around real hospitality moments: finding a place, joining an event, understanding a venue, paying locally, requesting support, and carrying a digital access pass.

The product should not be limited to hotel-room guests. It should model access, context, and venue relationships.

\begin{kazehighlightbox}
\textbf{Core principle:} Kaze should model access, not only accommodation. A user may be a hotel guest, event attendee, day-pass visitor, VIP invitee, conference participant, or public venue browser.
\end{kazehighlightbox}

\section{Target Users}

\begin{longtable}{p{0.30\textwidth} p{0.62\textwidth}}
\textbf{User type} & \textbf{Needs} \\
\midrule
Hotel guest & See stay schedule, request services, understand included or paid services, use maps, and access relevant hotel areas. \\
Wedding guest & Open an invitation code, view event information, receive updates, and use a Kaze Pass when entry is controlled. \\
Conference attendee & Find rooms, view schedule, access restricted event zones, and receive practical event details. \\
Venue shopper & Browse venues, compare capacity, pricing, location, and services before authentication. \\
Organizer & Create or manage event invitations, coordinate guests, request venue services, and connect bookings to access. \\
Venue or hotel operator & Publish spaces, manage booking interest, offer add-ons, and reduce phone-call friction for guests. \\
Local service provider & Offer services such as catering, styling, media, transport, cleaning, or insurance around venue bookings. \\
\end{longtable}

\section{Functional Scope}

\subsection{Public Discovery}

Kaze should allow users to explore useful information before forcing login. This is important for people who are comparing venues, apartments, event spaces, or local services.

\begin{itemize}
    \item Browse wedding venues and conference rooms.
    \item View basic venue details such as name, location, capacity, pricing, amenities, and service availability.
    \item Later support apartments and local stay alternatives.
    \item Provide a clear path from browsing into reservation, inquiry, invitation, or contact flows.
\end{itemize}

\subsection{Invitation and Code Entry}

Kaze should support invitation-based access and short-code entry.

\begin{itemize}
    \item Users can enter a short event or invitation code.
    \item Organizers can create an event and invite people, especially by phone number.
    \item Invitation codes should be generated only after the invitation or event has been confirmed and saved.
    \item Invited people should see relevant event details, updates, and pass state.
    \item Kaze Pass generation should depend on confirmation, approval, payment, or event rules when required.
\end{itemize}

\subsection{Venue Reservation Draft Flow}

Kaze should support reservation draft flows for venue categories where local booking is practical.

\begin{itemize}
    \item Wedding venue reservation request.
    \item Conference room reservation request.
    \item Amenity or day-pass booking direction.
    \item Add-on service selection during reservation.
    \item Payment preference capture using Rwanda-relevant methods.
\end{itemize}

\subsection{Stay and Hotel Services}

For hotel-linked users, Kaze should support stay details and service requests without requiring unnecessary input from the guest.

\begin{itemize}
    \item Show stay schedules, reserved services, and relevant venue locations.
    \item Distinguish included, free, paid, reserved, and optional services using simple language.
    \item Allow service requests such as towels, water, in-room dining, laundry pickup, concierge support, housekeeping, and custom requests.
    \item Avoid editable room number fields when the system already knows the user's room.
    \item Use notes or text areas where user context is useful, such as special instructions.
\end{itemize}

\section{Venue Map Requirements}

Kaze maps should be reusable across hotels, conference venues, wedding venues, stadiums, public buildings, and other structured spaces.

\begin{longtable}{p{0.30\textwidth} p{0.62\textwidth}}
\textbf{Map concept} & \textbf{Requirement} \\
\midrule
Site & A complete venue or property, such as a hotel, stadium, or public building. \\
Building or structure & A physical structure inside a site. \\
Floor or level & A navigable level within a structure. \\
Room or hall & A mapped space that can be reserved, visited, restricted, or described. \\
Zone & A defined area such as public, VIP, staff-only, event-only, or pass-required. \\
Seat or section & A future element for stadiums, conference layouts, and wedding seating. \\
Fixed object & Furniture or infrastructure that cannot be moved. \\
Movable layout object & Tables, chairs, aisles, and layout templates that can change per event. \\
Amenity or point of interest & Pool, gym, lounge, restaurant, restroom, gate, reception, or service point. \\
\end{longtable}

\subsection{Access-Aware Rendering}

Maps should support access rules, not only geometry.

\begin{itemize}
    \item Public areas render normally.
    \item Restricted areas can be marked or visually reduced.
    \item Hidden areas can be omitted entirely when the user should not know they exist.
    \item Access rules should eventually come from backend entitlements, not hardcoded local samples.
    \item Access can depend on guest type, event role, VIP level, venue policy, time window, reservation state, payment state, and Kaze Pass entitlements.
\end{itemize}

\section{Layout Planning and Event Styling}

Kaze should separate two different ideas that are often confused.

\begin{longtable}{p{0.30\textwidth} p{0.62\textwidth}}
\textbf{Term} & \textbf{Meaning} \\
\midrule
Layout planning & How chairs, tables, aisles, seating zones, fixed objects, and open spaces are arranged. This affects capacity, movement, and venue use. \\
Event styling & Flowers, lights, backdrops, table styling, theme decoration, and presentation services. This affects event look and service pricing. \\
\end{longtable}

For wedding venues and conference rooms, layout planning should eventually support:

\begin{itemize}
    \item round-table layouts;
    \item rectangular dining layouts;
    \item classroom layouts;
    \item theater seating;
    \item boardroom setups;
    \item standing reception zones;
    \item attendee-count suggestions;
    \item operator approval for layout choices.
\end{itemize}

\section{Payments and Commerce}

Kaze should support Rwanda-relevant payment and settlement flows.

\begin{itemize}
    \item MTN MoMo.
    \item Airtel Money.
    \item BK / RSwitch compatible flows.
    \item Card payments where available.
    \item Cash as a visible option where appropriate, with clear confirmation or validation flow.
    \item Deposits, balance collection, add-on service payments, and event-entry payments tied to access state.
\end{itemize}

\section{Business and Revenue Requirements}

Kaze should create value beyond simple venue listing so users and venues have fewer reasons to bypass the platform.

Potential revenue directions include:

\begin{itemize}
    \item venue discovery and booking commissions;
    \item partner subscription or SaaS fees for hotels and venues;
    \item add-on service marketplace fees;
    \item custom branded venue or hotel apps;
    \item event access and Kaze Pass infrastructure;
    \item map data and venue-layout tooling for other industries;
    \item lower-commission local alternatives to large international booking platforms where commercially useful.
\end{itemize}

\section{Non-Functional Requirements}

\begin{longtable}{p{0.30\textwidth} p{0.62\textwidth}}
\textbf{Area} & \textbf{Requirement} \\
\midrule
User experience & Premium, calm, fast, architectural, and easy to understand without long reading. \\
Responsiveness & Mobile-first, with tablet, foldable, and larger-screen layouts where needed. \\
Localization & English first, with future support for more languages. User-facing text should avoid developer-only wording. \\
Security & Private user data, access entitlements, event invitations, payment state, and partner data must be protected. \\
Privacy & The app should explain what data is collected and provide user controls where legally or product-wise appropriate. \\
Maintainability & Large UI files should be split into logical components and packages. Hardcoded demo data should move toward APIs and database-backed sources. \\
Extensibility & Venue maps, invitation themes, event types, and request flows should be data-driven where possible. \\
Notifications & Notification permission and notification delivery should be implemented only where the user benefits from timely updates. \\
\end{longtable}

\section{Backend and Data Requirements}

Kaze should move away from hardcoded data toward API and database-backed flows.

Important backend domains:

\begin{itemize}
    \item users and identities;
    \item venues, hotels, spaces, and amenities;
    \item maps, floors, rooms, zones, and access rules;
    \item reservations and booking drafts;
    \item events and invitation codes;
    \item Kaze Pass entitlements;
    \item service requests and request history;
    \item payment methods and payment state;
    \item notifications and realtime updates.
\end{itemize}

\begin{kazecalloutbox}[Architecture Direction]
Ktor is appropriate for the Kaze backend because it fits the Kotlin ecosystem and can support APIs, authentication, database access, realtime updates, and AI integrations when needed.
\end{kazecalloutbox}

\section{Future AI Scope}

AI can be added later where it improves the user experience, but it should not replace deterministic product flows.

Useful AI candidates:

\begin{itemize}
    \item answer practical guest questions such as whether the kitchen is open;
    \item explain venue rules or event details in simple language;
    \item help organizers draft invitations;
    \item summarize venue options based on budget, capacity, and location;
    \item support multilingual guest assistance.
\end{itemize}

AI responses should use trusted venue data and should not invent operational facts.

\section{Open Questions}

\begin{itemize}
    \item Which venue category should be monetized first: wedding venues, conference rooms, or hotel services?
    \item Which payment provider should be integrated first for production?
    \item Should public browsing require account creation before requesting a reservation?
    \item What approval rules are needed when one person creates an invitation for another organizer's event?
    \item What minimum legal documents are required before production launch?
    \item Which Kaze Pass formats should be prioritized: in-app pass, QR code, Apple Wallet, or Google Wallet?
\end{itemize}

\section{Contact}

\KazeContactBlock

\end{document}
